\documentclass{report}
\usepackage{amsmath,amssymb}
\setlength{\parindent}{0mm}
\setlength{\parskip}{1em}
\begin{document}
\begin{center}
\rule{6in}{1pt} \
{\large Erik G. Boman \\
{\bf Partitioning for Hybrid Direct-Iterative Preconditioners and Solvers}}

Sandia National Labs \\ PO Box 5800 \\ Albuquerque \\ NM 87185-1318
\\
{\tt egboman@sandia.gov}\\
Sivasankaran Rajamanickam\\
Jeremie Gaidamour\end{center}

Hybrid direct-iterative methods are attractive since they attempt
to combine ``the best of two worlds''. Several such solvers have been
developed recently. We focus on ShyLU and HIPS in this talk.
First we describe ShyLU, which will become a new Trilinos package.
We use a Schur complement method where subproblems are solved using a
direct solver while the Schur complement is solved iteratively.
We employ a block decomposition, which is formed in a fully automatic way
using state-of-the-art partitioning and ordering algorithms in
the Zoltan and Scotch software packages. We discuss the choice of
partitioning strategy, and in particular, narrow vs. wide separators,
and study the impact on the Schur complement system.

We have two main goals with ShyLU: First, we want a fast and robust
solver/preconditioner for multicore compute nodes. Second, due to its
modular design, ShyLU can serve as an research platform to compare
algorithms. For example, ShyLU supports two different approaches
to forming a sparse preconditioner for the Schur complement:
dropping and probing.

We demonstrate the robustness and effectiveness of ShyLU
for problems from a variety of applications. We have found
it is particularly effective on circuit simulation problems where
many iterative methods fail.

Finally, we compare ShyLU to HIPS, another parallel hybrid solver.
HIPS uses a hierarchical decomposition since it was designed for
large-scale problems, while ShyLU is intended to be used in conjunction
with domain decomposition or multigrid for very large problems.
We explain differences in the design and algorithms, and compare
their performance.


\end{document}
